%% kapa: exam=1 round=26 year=2015 period=2 form=A subject=appraisal_law subject_ko=부동산관계법규 seq=9 number=9
%% source: Q-Net 감정평가사 기출문제 (cst00309) artlSeq=5138001 file=extracted/1008/5138001/제26회_감정평가사 1차_기출/제26회_감정평가사 1차_기출/2교시_감평_A형.pdf
%% conversion: kordoc 4.24.51 -> tools/gen1.py (auto); review=proofread-sonnet; flags=none
\begin{question}{9}
국토의 계획 및 이용에 관한 법령상 도시ㆍ군계획시설사업 등에 관한 설명으로 옳은 것은?
\begin{choices}
  \item 도시ㆍ군계획시설사업의 시행자가 실시계획의 인가를 받고자 하는 경우 국토교통부장관이 지정한 시행자는 시ㆍ도지사의 인가를 받아야 한다.
  \item 준공검사를 받은 후에 해당 도시ㆍ군계획시설사업에 대하여 사업명칭을 변경하기 위하여 실시계획을 작성하는 경우에도 국토교통부장관의 인가를 받아야 한다.
  \item 도시ㆍ군계획시설사업에 관한 실시계획에는 사업의 종류 및 명칭은 포함되어야 하지만, 사업의 면적 또는 규모는 반드시 포함되어야 하는 것은 아니다.
  \item 행정청이 아닌 도시ㆍ군계획시설사업 시행자의 처분에 대해서는 그 시행자에게 행정심판을 제기하여야 한다.
  \item 도시ㆍ군계획시설사업을 분할시행하는 때에는 분할된 지역별로 실시계획을 작성할 수 있다.
\end{choices}
\end{question}
\answer{5}
